\chapter{Justificación}
\label{ch:justificacion}

La relevancia de esta investigación radica en la necesidad imperativa de proteger los activos digitales universitarios mediante enfoques de ingeniería formal y estándares de la industria \citep{iso202227001,nist2018cybersecurity}.

\section{Justificación Técnica}
Desde la perspectiva de la ingeniería de sistemas y la ciberseguridad, el modelado formal de amenazas es una de las prácticas más costo-efectivas para elevar la postura de seguridad \citep{howard2003writing,shostack2014threat}. La identificación de vectores de ataque en las capas de aplicación, transporte y almacenamiento permite diseñar controles de seguridad en profundidad (\textit{Defense in Depth}), implementando políticas de autenticación multifactor (MFA), sanitización estricta de entradas y segmentación de redes, reduciendo drásticamente la probabilidad de compromisos críticos frente a los riesgos del OWASP Top 10 \citep{owasp2021top10}.

\section{Justificación Académica}
Este estudio contribuye al cuerpo de conocimiento técnico de la Escuela Profesional de Ingeniería de Sistemas de la UNAP, proporcionando un marco metodológico reproducible que demuestra cómo integrar estándares industriales de primer nivel (Microsoft Threat Modeling \citep{kohnfelder1999threats} y Lockheed Martin Cyber Kill Chain \citep{hutchins2011intelligence}) en proyectos de evaluación real y auditoría académica.

\section{Justificación Institucional y Social}
La salvaguarda de la integridad de las calificaciones, los registros contables y los datos personales de la comunidad universitaria de la UNAP es una obligación legal y ética enmarcada en la Ley N° 29733 de Protección de Datos Personales del Perú \citep{peru2011ley29733}. Este trabajo provee a las autoridades y a la Oficina de Tecnologías de la Información (OTI-UNAP) de un diagnóstico técnico claro para proteger el patrimonio informático institucional.
